\chapter*{Acknowledgments}
\addcontentsline{toc}{chapter}{Acknowledgments}

We wish to express our sincere gratitude to the individuals and organisations
whose support made this capstone project possible.

First and foremost, we thank the course instructors and teaching assistants of
\textbf{AT82.9002 Data Engineering and MLOps} at the Asian Institute of
Technology for their guidance throughout the semester. The continuous feedback
on our proposal and progress report shaped the scope, rigour, and engineering
discipline of the final CAAS system. Specific feedback on drift-monitoring
policy design, cost realism, and contribution framing was instrumental in
strengthening the final report.

We gratefully acknowledge the data providers whose open access made this work
feasible: the \textbf{Pollution Control Department (PCD) of Thailand} for
publishing historical PM2.5 records through the air4thai portal;
\textbf{Open-Meteo} for providing free ERA5-quality meteorological reanalysis
through a well-documented API; and \textbf{NASA FIRMS} for making
near-real-time satellite fire hotspot data available under a permissive
research licence. Without these three sources, a self-contained, reproducible
end-to-end pipeline would not have been possible at the student scale.

We thank the open-source maintainers of the libraries and services on which
CAAS is built --- LightGBM, XGBoost, TensorFlow, scikit-learn, Optuna, MLflow,
Evidently, FastAPI, Streamlit, Terraform, Docker, and GitHub Actions ---
whose collective work compresses the distance between a research prototype
and a production-grade system into a matter of weeks of student effort.

Finally, we thank our families, classmates, and peers at AIT for their
patience and encouragement during long debugging sessions, cloud
provisioning failures, and the inevitable late-night rebuilds that
every real engineering project entails.

\vspace{0.5cm}

\noindent Supanut Kompayak (st126055) \\
\noindent Shuvam Shrestha (st125975) \\
\noindent Asian Institute of Technology, April 2026
